\documentclass{article}
\usepackage{graphicx}
\usepackage[margin=1.5cm]{geometry}
\usepackage{amsmath}
\usepackage{enumitem}

\setlist{nosep,leftmargin=*}

\begin{document}
\twocolumn

\title{Chapter 5 Warm-Up: Python \texttt{for} and \texttt{while} Loops}
\author{Prof. Jordan C. Hanson}
\date{}

\maketitle

\section{Student Data}

\begin{enumerate}

\item A list can contain tuples representing individual students.
Each tuple below contains a student's name, student ID, and math
placement score:

\begin{verbatim}
students = [
 ("Maya", 20681001, 5),
 ("Noah", 20681002, 1),
 ("Olivia", 20681003, 3),
 ("Liam", 20681004, 4),
 ("Sophia", 20681005, 2)
]
\end{verbatim}

Each student therefore has the form

\begin{verbatim}
(name, student_id, score)
\end{verbatim}

Try

\begin{verbatim}
print(students[0])
print(students[0][0])
print(students[0][1])
print(students[0][2])
\end{verbatim}

Identify which expression produces the student's name,
ID, and math score.

\end{enumerate}


\section{\texttt{for} Loops}

\begin{enumerate}

\item A \texttt{for} loop performs the same operation on
each element of a collection.  Enter

\begin{verbatim}
for student in students:
    print(student)
\end{verbatim}

The variable \texttt{student} refers to one tuple during
each pass through the loop.

\begin{itemize}

\item Modify the loop to extract all three values:

\begin{verbatim}
for student in students:
    name = student[0]
    student_id = student[1]
    score = student[2]
    print(name, student_id, score)
\end{verbatim}

\item Add conditional logic inside the loop so that each
student's recommended course is printed:

\begin{verbatim}
if score >= 4:
    print("Calculus")
elif score >= 2:
    print("Trigonometry")
else:
    print("Algebra")
\end{verbatim}

\end{itemize}

\end{enumerate}


\section{\texttt{while} Loops}

\begin{enumerate}

\item A \texttt{while} loop repeats while a Boolean
condition remains \texttt{True}.  An index can be used
to select one student at a time:

\begin{verbatim}
i = 0

while i < 5:
    print(students[i])
    i = i + 1
\end{verbatim}

\begin{itemize}

\item What value does \texttt{i} have during the first
pass through the loop?  During the last pass?

\item What would happen if

\begin{verbatim}
i = i + 1
\end{verbatim}

were omitted?

\item Modify the loop so that it prints only each
student's name and student ID.

\end{itemize}

\end{enumerate}


\section{Warm-Up Exercise}

\begin{enumerate}

\item Create the following initially empty course structures:

\begin{verbatim}
algebra = ["Algebra", []]
trig = ["Trigonometry", []]
calculus = ["Calculus", []]

math_courses = (
    algebra,
    trig,
    calculus
)
\end{verbatim}

Use a \texttt{for} loop to examine every student in
\texttt{students}.  Use an \texttt{if}-\texttt{elif}-\texttt{else} statement to determine the correct course. Instead of adding only the student's name, add a tuple containing the student's name and ID:

\begin{verbatim}
(name, student_id)
\end{verbatim}

For example,

\begin{verbatim}
calculus[1].append(
    (name, student_id)
)
\end{verbatim}

When the loop is finished, print

\begin{verbatim}
print(math_courses)
\end{verbatim}

and verify that

\begin{itemize}
\item scores 0--1 are placed in Algebra,
\item scores 2--3 are placed in Trigonometry,
\item scores 4--5 are placed in Calculus, and
\item every roster contains both student names and IDs.
\end{itemize}

\end{enumerate}

\end{document}